% GENERATED FILE. Edit canonical lessons, then run npm run generate:books.
% canonical-source-hash: fnv1a64:1e7984e81aea3983
% canonical-lessons: ZH-C189-hao, ZH-W189-hao-number, ZH-C189-biaoge, ZH-C189-tianxie, ZH-C189-qianming, ZH-C189-shenfenzheng

\chapter{Forms and Dates}
\label{ch:zh-forms-and-dates}
\hlchaptermodality{189}

\begin{chapteropening}
\textbf{By the end of this chapter:} \emph{I can say the date, ask for a form, fill it in, sign it and show my identity card.}

Complete the last lesson of chapter 189: I can say the date, ask for a form, fill it in, sign it and show my identity card.
\end{chapteropening}

\section[\texorpdfstring{\zh{号} (hào)}{号 (hào)}]{\zh{号} (hào) — day of the month; number (in a series)}

\label{lesson:ZH-C189-hao}



\begin{quote}
Before the new one: say the Mandarin for to spend less, then the Mandarin for enough.
\end{quote}

\subsection*{You'll want to know: \zh{号}}
\textbf{\zh{号}} — \emph{hào} — \textquotedblleft{}day of the month; number (in a series)\textquotedblright{}.

\begin{quote}
\textbf{\zh{今天是十月五号}} \emph{jīntiān shì shíyuè wǔ hào} — Today is the fifth of October.
\end{quote}

The new character, in pieces: \zh{号} — five strokes: \zh{口} on top, then a long bar, and under it a stroke that goes down, turns right, turns down again and hooks to the left. The next lesson puts it in your hand.

\begin{cousinweb}
\textbf{\zh{号}} (\emph{hào}) is inherited from Middle Chinese \emph{hawH} (Guangyun, in Baxter's transcription), glossed 'to call', 'a command' and 'a title'. Numbering days and things is not given there, so how it came to that sense is not settled.
\end{cousinweb}

\begin{grammarlens}[title={What you've built}]
One more word for reading notices and messages.
\end{grammarlens}

\subsection*{Your turn}
\begin{itemize}
\raggedright
  \item \emph{Say it:} \emph{hào}
  \item \emph{Say it:} \emph{hào}, once more
  \item \emph{Say it:} \emph{jīntiān shì shíyuè wǔ hào} — Today is the fifth of October.
  \item \emph{Recall:} say the Mandarin for to spend less, then the Mandarin for enough, then say \emph{hào} again
\end{itemize}

\begin{tcolorbox}[breakable,colback=teal!4,colframe=teal!35!black,title={Before you move on}]
What does \zh{号} mean? (\textquotedblleft{}Day of the month; number (in a series)\textquotedblright{}.) Say it once more.
\end{tcolorbox}
\section[\texorpdfstring{\zh{号} (hào)}{号 (hào)}]{\zh{号} — \zh{口} on top, then a bar and a hook}

\label{lesson:ZH-W189-hao-number}



\begin{quote}
The word first, then the character inside it.

\begin{itemize}
\raggedright
  \item \emph{Recall:} say \emph{hào} — \textbf{R1}, one lesson back, before it decays
\end{itemize}

\begin{quote}
\zh{月} \zh{口}
\end{quote}

\emph{Read it:} each one aloud, with nothing said first — \emph{yuè}, \emph{kǒu}
\end{quote}

\subsection*{Script — \zh{口} on top, then a bar and a hook}
\hlblockfigure{\detokenize{figures/ZH-W189-hao-number-filmstrip.pdf}}{How \zh{号} is written, stroke by stroke}

\begin{quote}
\zh{号}
\end{quote}

\textbf{Two pieces, top before bottom.}

\textbf{The top piece, \zh{口}, three strokes,} the same three as in \zh{叫} and \zh{吃}: a short vertical, a stroke that runs right and turns down, then the bar that closes the foot.

\textbf{The bottom piece, two strokes.} A long bar across the middle, wider than \zh{口}. Then, from under the bar, a short stroke down that turns right, turns down again and hooks to the left at its foot.

3 + 2 = five.

\subsection*{Your turn}
\begin{itemize}
\raggedright
  \item \emph{Write it:} \zh{号}, all five, with the model in front of you
  \item \emph{Cover:} write it again from memory, and count the pieces, not the strokes
\end{itemize}

\begin{tcolorbox}[breakable,colback=teal!4,colframe=teal!35!black,title={Before you move on}]
How many strokes in \zh{号}? (\textbf{Five}.) Which piece comes first? (\textbf{\zh{口}, on top}.) Where does the last stroke hook? (\textbf{To the left, at its foot}.)

Source: \href{https://raw.githubusercontent.com/chanind/hanzi-writer-data/68d10a4b21150cae5e1ebbd223eed289cf32d90c/data/\%E5\%8F\%B7.json}{Hanzi Writer Data \zh{号}, pinned stroke paths and medians}.
\end{tcolorbox}
\section[biǎogé]{biǎogé — a form to fill in}

\label{lesson:ZH-C189-biaoge}



\begin{quote}
Before the new one: say the Mandarin for enough, then the Mandarin for day of the month.
\end{quote}

\subsection*{You'll want to know: biǎogé}
\emph{biǎogé} — \textquotedblleft{}a form to fill in\textquotedblright{}.

\begin{quote}
\emph{qǐng gěi wǒ yì zhāng biǎogé} — Could I have a form, please?
\end{quote}

Its characters come later; for now it is a word for the ear.

\begin{cousinweb}
\emph{biǎogé} is built from \emph{biǎo} and \emph{gé}, both inherited from Middle Chinese, \emph{pjewX} and \emph{kaek} (Guangyun, in Baxter's transcription). The dictionary glosses the first 'to make clear' and the second 'a pattern' and 'a measure'; how the pair came to mean a form to fill in is not settled.
\end{cousinweb}

\begin{grammarlens}[title={What you've built}]
One more word for reading notices and messages.
\end{grammarlens}

\subsection*{Your turn}
\begin{itemize}
\raggedright
  \item \emph{Say it:} \emph{biǎogé}
  \item \emph{Say it:} \emph{biǎogé}, once more
  \item \emph{Say it:} \emph{qǐng gěi wǒ yì zhāng biǎogé} — Could I have a form, please?
  \item \emph{Recall:} say the Mandarin for enough, then the Mandarin for day of the month, then say \emph{biǎogé} again
\end{itemize}

\begin{tcolorbox}[breakable,colback=teal!4,colframe=teal!35!black,title={Before you move on}]
What does biǎogé mean? (\textquotedblleft{}A form to fill in\textquotedblright{}.) Say it once more.
\end{tcolorbox}
\section[tiánxiě]{tiánxiě — to fill in (a form)}

\label{lesson:ZH-C189-tianxie}



\begin{quote}
Before the new one: say the Mandarin for day of the month, then the Mandarin for a form to fill in.
\end{quote}

\subsection*{You'll want to know: tiánxiě}
\emph{tiánxiě} — \textquotedblleft{}to fill in (a form)\textquotedblright{}.

\begin{quote}
\emph{qǐng tiánxiě zhè zhāng biǎogé} — Please fill in this form.
\end{quote}

Its characters come later; for now it is a word for the ear.

\begin{cousinweb}
\emph{tiánxiě} is built from \emph{tián} 'to fill up' and \emph{xiě} 'to write', both inherited from Middle Chinese, \emph{den} and \emph{sjaeX} (Guangyun, in Baxter's transcription). The dictionary glosses the first 'to stop up' and 'full', and the second, among other senses, 'to copy out a book'.
\end{cousinweb}

\begin{grammarlens}[title={What you've built}]
One more word for reading notices and messages.
\end{grammarlens}

\subsection*{Your turn}
\begin{itemize}
\raggedright
  \item \emph{Say it:} \emph{tiánxiě}
  \item \emph{Say it:} \emph{tiánxiě}, once more
  \item \emph{Say it:} \emph{qǐng tiánxiě zhè zhāng biǎogé} — Please fill in this form.
  \item \emph{Recall:} say the Mandarin for day of the month, then the Mandarin for a form to fill in, then say \emph{tiánxiě} again
\end{itemize}

\begin{tcolorbox}[breakable,colback=teal!4,colframe=teal!35!black,title={Before you move on}]
What does tiánxiě mean? (\textquotedblleft{}To fill in (a form)\textquotedblright{}.) Say it once more.
\end{tcolorbox}
\section[qiānmíng]{qiānmíng — to sign your name}

\label{lesson:ZH-C189-qianming}



\begin{quote}
Before the new one: say the Mandarin for a form to fill in, then the Mandarin for to fill in.
\end{quote}

\subsection*{You'll want to know: qiānmíng}
\emph{qiānmíng} — \textquotedblleft{}to sign your name\textquotedblright{}.

\begin{quote}
\emph{nǐ qiānmíng le ma} — Have you signed it?
\end{quote}

Its characters come later; for now it is a word for the ear.

\begin{cousinweb}
\emph{qiānmíng} is built from \emph{qiān} 'to sign' and \emph{míng} 'a name'. The second is inherited from Middle Chinese \emph{mjieng} (Guangyun, in Baxter's transcription); no dictionary we use records an older form of the first, so its history is not settled.
\end{cousinweb}

\begin{grammarlens}[title={What you've built}]
One more word for reading notices and messages.
\end{grammarlens}

\subsection*{Your turn}
\begin{itemize}
\raggedright
  \item \emph{Say it:} \emph{qiānmíng}
  \item \emph{Say it:} \emph{qiānmíng}, once more
  \item \emph{Say it:} \emph{nǐ qiānmíng le ma} — Have you signed it?
  \item \emph{Recall:} say the Mandarin for a form to fill in, then the Mandarin for to fill in, then say \emph{qiānmíng} again
\end{itemize}

\begin{tcolorbox}[breakable,colback=teal!4,colframe=teal!35!black,title={Before you move on}]
What does qiānmíng mean? (\textquotedblleft{}To sign your name\textquotedblright{}.) Say it once more.
\end{tcolorbox}
\section[shēnfènzhèng]{shēnfènzhèng — identity card}

\label{lesson:ZH-C189-shenfenzheng}



\begin{quote}
Before the new one: say the Mandarin for to fill in, then the Mandarin for to sign your name.
\end{quote}

\subsection*{You'll want to know: shēnfènzhèng}
\emph{shēnfènzhèng} — \textquotedblleft{}identity card\textquotedblright{}.

\begin{quote}
\emph{nǐ dài shēnfènzhèng le ma} — Have you brought your ID card?
\end{quote}

Its characters come later; for now it is a word for the ear.

\begin{cousinweb}
\emph{shēnfènzhèng} is built from \emph{shēn} 'oneself', \emph{fèn} 'a share, a standing' and \emph{zhèng} 'proof'. The first and the last are inherited from Middle Chinese, \emph{syin} and \emph{tsyingH} (Guangyun, in Baxter's transcription), glossed 'oneself' and 'to verify'; the dictionary gives the middle character another sound and sense, so its history is not settled.
\end{cousinweb}

\begin{grammarlens}[title={What you've built}]
One more word for reading notices and messages.
\end{grammarlens}

\subsection*{Your turn}
\begin{itemize}
\raggedright
  \item \emph{Say it:} \emph{shēnfènzhèng}
  \item \emph{Say it:} \emph{shēnfènzhèng}, once more
  \item \emph{Say it:} \emph{nǐ dài shēnfènzhèng le ma} — Have you brought your ID card?
  \item \emph{Recall:} say the Mandarin for to fill in, then the Mandarin for to sign your name, then say \emph{shēnfènzhèng} again
\end{itemize}

\begin{tcolorbox}[breakable,colback=teal!4,colframe=teal!35!black,title={Before you move on}]
What does shēnfènzhèng mean? (\textquotedblleft{}Identity card\textquotedblright{}.) Say it once more.
\end{tcolorbox}
